\section{The finitely presented residually finite 2-group}
\label{sec:construction}

Stable detection and the fixed-rank injectivity just proved now yield
finite quotients separating the group we will use. The ring $R$
throughout this section is the algebra of
Proposition~\ref{prop:ring-input}, with the hierarchy of
Proposition~\ref{prop:hierarchy-input}. Its finite degree quotients
will give actual finite groups separating the elements of
$\Gamma=\St_{12}(R)$.

\begin{proposition}\label{prop:finite-quotients}
For every $d\ge1$, the group $\St_{12}(R/R_{\ge d})$ is finite.
The homomorphisms
\[
 \St_{12}(R)\longrightarrow\St_{12}(R/R_{\ge d}),\qquad d\ge1,
\]
separate the elements of $\St_{12}(R)$.
\end{proposition}

\begin{proof}
Put $B_d=R/R_{\ge d}$. This ring is finite: it is the direct sum of
finitely many finite-dimensional vector spaces over $\F_2$.
Moreover $B_d=\F_2\oplus I_d$ with $I_d^d=0$, so every finite
square matrix over $I_d$ is nilpotent. The centrality assertion in
Proposition~\ref{prop:finite-rank-input} therefore applies to $B_d$.
The kernel of
\[
 \St_{12}(B_d)\longrightarrow\E_{12}(B_d)
\]
is central, and the matrix group on the right is finite. Hence the
center of $\St_{12}(B_d)$ has finite index. Schur's theorem states
that a group with center of finite index has finite derived subgroup
\cite{Schur1904}; see also \cite[Theorem~10.1.4]{Robinson1996}. Since $\St_{12}(B_d)$ is perfect, it is itself
finite.

Let $g\ne1$ in $\St_{12}(R)$. Theorem~\ref{thm:stability} makes
its image in $\St(R)$ nonidentity, and
Theorem~\ref{thm:detection} detects that image in $\St(B_d)$ for
some $d$. The natural square
\[
 \begin{tikzpicture}[baseline=(current bounding box.center),
                    node distance=1.4cm and 2.7cm]
 \node (a) {$\St_{12}(R)$};
 \node (b) [right=of a] {$\St(R)$};
 \node (c) [below=of a] {$\St_{12}(B_d)$};
 \node (d) [right=of c] {$\St(B_d)$};
 \draw[->] (a)--(b);
 \draw[->] (a)--(c);
 \draw[->] (b)--(d);
 \draw[->] (c)--(d);
 \end{tikzpicture}
\]
commutes on each root generator. Thus the image of $g$ in
$\St_{12}(B_d)$ is nonidentity. No injectivity assertion for the
bottom arrow is needed. This proves separation by the finite groups
just constructed.
\end{proof}

\begin{proof}[Proof of Theorem~\ref{thm:main}]
The algebra $R$ is infinite and finitely presented as a unital
$\F_2$-algebra. Accordingly $\Gamma=\St_{12}(R)$ has an ordinary
finite presentation by \cite[Theorem~4.1]{Companion}; this is also a
case of \cite[Theorem~3]{KrsticMcCool}. The group $\Gamma$ is
infinite, since its matrix image contains distinct matrices
$e_{12}(a)$ for all $a\in R$. It is residually finite by
Proposition~\ref{prop:finite-quotients}.

Define
\[
 G=\ker\bigl(\Gamma\longrightarrow\St_{12}(\F_2)\bigr),
\]
where the map is induced by the degree-zero projection of $R$.
The case $d=1$ of Proposition~\ref{prop:finite-quotients} shows
that this target is finite. Thus $G$ has finite index in $\Gamma$.
It is infinite and has an ordinary finite presentation by the
Reidemeister--Schreier theorem \cite[Chapter~II, Section~4]{LyndonSchupp}. Restricting the
separating finite-quotient maps of $\Gamma$ proves residual finiteness
of $G$.

Fix $g\in G$. Its matrix image is $\Id+A$ for some
$A\in M_{12}(I)$. Proposition~\ref{prop:ring-input} gives an integer
$r\ge1$ with $A^r=0$. Choose $a\ge0$ with $2^a\ge r$.
In characteristic two,
\[
 (\Id+A)^{2^a}=\Id+A^{2^a}=\Id.
\]
Consequently $g^{2^a}\in J_{12}(R)$, and its reduction to
$\St_{12}(\F_2)$ is still trivial. Its stable image belongs to
$J(R)$ and reduces to $1$ in $\St(\F_2)$, so
Proposition~\ref{prop:relative-torsion} kills that image by a
2-power $2^b$. The injectivity in Theorem~\ref{thm:stability}
then yields $g^{2^{a+b}}=1$ in $\Gamma$, hence in $G$.
This proves every assertion of Theorem~\ref{thm:main}.
\end{proof}
